\chapter{كيف تتعلم الآلة الحركة}
\chaptermark{كيف تتعلم الآلة الحركة}
\label{sec:motorlearning}
\begin{chaptergoals}
\item لماذا تحتاج الحركة إلى معلّم ثالث، سلك خطأ إلى كل مخرج، وما ثمنه؛
\item كيف تعطي طبقة توسيع هائلة من \nt{GLUT} وسلك خطأ لكل مخرج قاعدة المربعات الصغرى؛
\item كيف تصير الدارة بعد التعلّم نموذجًا داخليًا، مرشحًا تكيّفيًا يتنبأ بالجسم؛
\item لماذا تفسر عمليتا تعلّم، سريعة وبطيئة، الأثر اللاحق والاستعادة التلقائية.
\end{chaptergoals}

\section{ثلاثة معلّمين}

رأينا في الفصل~\ref{sec:muscles} كيف تحرك الآلة العضلات، وتوقفنا عند فرضية: الحركات السريعة الدقيقة تتطلب نموذجًا داخليًا للجسم يتنبأ بنتيجة الأمر~\cite{WolpertGhahramani2000}. فمن أين يأتي هذا النموذج؟ يجب أن يُتعلَّم، وأن يُتعلَّم على غير الطريقة التي تعلمنا بها حتى الآن.

مرّ في الكتاب معلّمان. الأول \emph{لا أحد}: قاعدة هيب (الفصل~\ref{sec:dofa}) تقوّي ما يحدث معًا كثيرًا، وتضغط المداخل. والثاني \emph{المكافأة}: \nt{DOPA} يبث إلى الجميع عددًا واحدًا، ”أفضل أو أسوأ مما كان متوقعًا“. أما الحركة فتحتاج إلى معلّم ثالث هو \emph{الخطأ}. أخطأت اليد الفنجان بسنتيمترين إلى اليسار: هذا ليس ”أسوأ“، بل متجه يقول لكل مخرج أين أخطأ وبكم. وكان المهندس سيسمي هذا التعلم الموجَّه.

الفرق بين المعلّم الثاني والثالث هو الفرق بين سلك واحد للجميع وسلك مستقل لكل مخرج. في القضية~\ref{prop:broadcastcost} كان التعلم بعدد مشترك واحد يبطؤ مع كل عصبون يؤثر ضجيجه في النتيجة. أما إذا تلقى كل مخرج $i$ خطأه الخاص $e_i$، فيمكن تغيير الأوزان بأبسط قاعدة:
\begin{empheq}[box=\keybox]{equation}
\frac{\dd w_{ij}}{\dd t} \;=\; -\eta\,e_i(t)\,x_j(t) .
\label{eq:lms}
\end{empheq}
هذه قاعدة المربعات الصغرى المعروفة منذ زمن في تقنية المرشحات التكيّفية~\cite{WidrowHoff1960}: يتغير الوزن بما يتناسب مع حاصل ضرب دخله في خطأ مخرجه، ويسير في المتوسط مع تدرج مربع الخطأ. ولا حاجة إلى ضجيج للبحث كما في الفصل~\ref{sec:dofa}: اتجاه التصحيح يأتي مباشرة عبر سلك الخطأ. ولا تنخفض السرعة مع عدد المخارج، لأن أخطاء الآخرين لا تصل إلى المشبك.

\begin{keyblock}
\begin{proposition}[ثلاثة معلّمين]
\label{prop:teachers}
قاعدة هيب تعلّم بلا معلّم ما هو متكرر؛ وإشارة \nt{DOPA} تعلّم بعدد واحد للشبكة كلها ما هو مجدٍ، وتبطؤ مع كل عصبون؛ وسلك الخطأ إلى كل مخرج يعلّم بالقاعدة~\eqref{eq:lms} ما هو دقيق، وسرعته لا تتوقف على عدد المخارج. وثمن الطريقة الثالثة سلك مستقل لكل مخرج، ومن يعرف الجواب الصحيح.
\end{proposition}
\end{keyblock}

\section{سلك الخطأ}

من يستطيع أن يعرف الجواب الصحيح للحركة؟ المستشعرات نفسها. إذا انحرفت اليد يسار الهدف أبلغت عن ذلك العينُ ومستشعرات التمدد من الفصل~\ref{sec:sensors}. والمطلوب دارة توصل هذا الخطأ إلى المخارج المعنية.

في الدماغ منطقة يبدو أنها مبنية لهذا بالضبط~\cite{Marr1969,Albus1971}. دارتها منتظمة على نحو غير مألوف، تكاد تكون بلورية، ومن السهل أن نتعرف فيها على القاعدة~\eqref{eq:lms} (الشكل~\ref{fig:errorwire}).

\emph{طبقة توسيع الدخل.} تصل الإشارات عن الأمر وعن حالة الجسم إلى طبقة هائلة من عصبونات \nt{GLUT} الصغيرة. عددها نحو سبعين مليارًا، أي أكثر مما في بقية الدماغ مجتمعة~\cite{Azevedo2009}. كل منها يمزج بضعة مداخل، فتنبسط الإشارة في متجه طويل جدًا. والمهندس يعرف هذه الحيلة: إذا رُفع بُعد الدخل، أمكن في الفضاء الجديد الحصول على أي علاقة مطلوبة تقريبًا بمجموع موزون بسيط~\cite{Cover1965}.

\emph{عصبونات المخرج.} يحسب المجموع الموزون نحو ثلاثين مليون عصبون مخرج كبير~\cite{Andersen1992cereb}، لكل منها نحو مئة ألف مدخل من طبقة التوسيع. وهي عصبونات \nt{GABA}: نتيجة التعلم تُنقل إلى الأمام لا بالإثارة بل بالتثبيط، كما في المنع من الفصل~\ref{sec:gaba}.

\emph{سلك الخطأ.} يتلقى كل عصبون مخرج مدخلًا آخر خاصًا: سلك \nt{GLUT} واحد بالضبط، يطلق نادرًا، نحو مرة في الثانية، لكنه يستدعي في كل مرة ومضة قوية في عصبون المخرج~\cite{Ito2001}. هذا هو $e_i$: احتمال الومضة يتوقف على اتجاه خطأ الحركة~\cite{Herzfeld2018}. وتزامن ومضة الخطأ مع نشاط المدخل يُضعف هذا المدخل، وهذا بالضبط الحد $-\eta\,e_i x_j$ في~\eqref{eq:lms}.

\begin{figure}[t]
\centering
\begin{tikzpicture}[>=Stealth,
  gran/.style={circle,draw,thick,fill=blue!6,minimum size=0.32cm,inner sep=0pt},
  outn/.style={circle,draw,very thick,fill=red!10,minimum size=1.1cm,inner sep=0pt,font=\scriptsize},
  inh/.style={-{Bar[width=7pt]},very thick,red!70!black}]
% inputs
\foreach \y/\lab in {1.6/{الأمر},0.4/{حالة الجسم}}{
  \draw[->,thick,blue!60!black] (-0.6,\y) node[left,font=\scriptsize,black] {\lab} -- (0.35,\y);}
% expansion layer
\foreach \i in {0,...,11}{\node[gran] (g\i) at ({0.8+0.28*mod(\i,2)},{-0.3+0.23*\i}) {};}
\foreach \i in {0,...,11}{\pgfmathsetmacro{\yy}{mod(\i,3)>0 ? 1.6 : 0.4}\draw[gray!60!black] (0.35,\yy) -- (g\i);}
\node[font=\scriptsize,align=center] at (0.95,-1.05) {طبقة التوسيع\\$\sim7\cdot10^{10}$ \nt{GLUT}};
% parallel wires to the output neuron
\foreach \i in {0,...,11}{\draw[->,gray!70!black] (g\i.east) -- (4.1,{1.0+0.07*(\i-5.5)});}
\node[outn] (P) at (4.6,1.0) {\nt{GABA}};
\node[font=\scriptsize,align=center,above=2pt] at (P.north) {عصبون المخرج\\$\sim10^5$ مدخل $x_j$، أوزان $w_{ij}$};
\draw[inh] (P) -- (6.8,1.0) node[right,font=\scriptsize,align=left] {تصحيح\\للحركة};
% error wire
\draw[->,very thick,red!70!black,densely dashed] (4.6,-1.4) node[below,font=\scriptsize,black,align=center] {سلك خطأ واحد $e_i$\\\nt{GLUT}، $\sim1$ مرة/ث} -- (P.south);
\end{tikzpicture}
\caption{دارة التعلم الموجَّه كما يبدو أن الدماغ ينفذها. تبسط طبقة هائلة من عصبونات \nt{GLUT} الأمرَ وحالةَ الجسم في متجه طويل $x_j$؛ ويجمعه عصبون \nt{GABA} المخرج بالأوزان $w_{ij}$. ويأتي سلك الخطأ الوحيد بـ$e_i$؛ وتزامن الخطأ مع مدخل نشط يُضعف وزن هذا المدخل، كما في القاعدة~\eqref{eq:lms}.}
\label{fig:errorwire}
\end{figure}

وهناك صعوبة نعرفها من الفصل~\ref{sec:dofa}: الخطأ يصل متأخرًا عن الدخل الذي قاد إليه. فالإخفاق يظهر بعد عشرات إلى مئات الملّي ثانية من الأمر. إذن يجب أن يتذكر المشبك أنه كان نشطًا، أي يلزم أثر، كما في القاعدة~\eqref{eq:threefactor}. وطول هذا الأثر، بحسب التجارب، مضبوط على تأخير التغذية الراجعة في كل مهمة: حيث يصل الخطأ بعد مئة ملّي ثانية يضعف أكثر ما يضعف المدخلُ الذي كان نشطًا قبله بمئة ملّي ثانية~\cite{Suvrathan2016}.

\begin{keyblock}
\begin{proposition}[سلك الخطأ]
\label{prop:errorwire}
في الدارة المبيّنة في الشكل~\ref{fig:errorwire} يتلقى كل عصبون مخرج سلك خطأ واحدًا بالضبط، وتزامن الخطأ مع المدخل يُضعف هذا المدخل. هذه قاعدة المربعات الصغرى~\eqref{eq:lms} مطبَّقة على دخل وُسِّع إلى بُعد هائل. بنية الدارة والإضعاف عند التزامن مقيسان؛ أما أن الدارة كلها تعمل تعلمًا موجَّهًا ففرضية رائدة.
\end{proposition}
\end{keyblock}

\section{النموذج الداخلي مرشحًا تكيّفيًا}

ماذا تتعلم دارة كهذه؟ أقرب جواب: التنبؤ. ليصل إلى الدخل الأمر $u(t)$ مارًّا بمجموعة من المرشحات بثوابت زمنية مختلفة، كما في طبقة التوسيع: $x_j(t)=(g_j*u)(t)$. والمخرج مجموعها الموزون، أي التنبؤ بما ستُبلغ عنه المستشعرات:
\begin{empheq}[box=\keybox]{equation}
\hat y(t) \;=\; \sum_j w_j\,(g_j*u)(t), \qquad e(t) \;=\; \hat y(t) - y(t) ,
\label{eq:adaptivefilter}
\end{empheq}
والأوزان تتعلم بحسب~\eqref{eq:lms}. هذا \emph{مرشح تكيّفي}: دارة تختار بنفسها دالة تحويلها لتكرر نظامًا مجهولًا~\cite{Fujita1982}. والنظام المجهول هنا هو الجسم نفسه بكتلته ومرونته وتأخيراته. والمرشح المتعلَّم هو النموذج الداخلي من الفصل~\ref{sec:muscles}: يقول ما ستُبلغ عنه المستشعرات دون انتظارها.

ونموذج كهذا مفيد مرتين. أولًا، يمكن تقديم تنبؤه بدل التغذية الراجعة المتأخرة، فلا يعود شرط الاستقرار~\eqref{eq:delaystable} يقيّد اليدين: يمكن التصحيح بسرعة، وفق النموذج. وثانيًا، الفرق بين المتنبَّأ به والمتلقَّى إشارة عن \emph{غير المتوقع}. فكل ما فعلته الآلة بنفسها تنبأ به النموذج وطرحه؛ ويبقى ما فعله العالم. ولهذا، بحسب إحدى الفرضيات، لا نستطيع أن ندغدغ أنفسنا~\cite{Blakemore1998}.

\section{حين يتغير العالم: السريع والبطيء}

لنختبر الدارة بتجربة سهلة الإعداد. يمدّ الإنسان يده نحو هدف، فيتغير العالم فجأة: مثلًا، تؤثر في اليد قوة جانبية. تخطئ الحركات الأولى، ثم يتناقص الخطأ. وإذا أُزيلت القوة، أخطأت اليد أولًا في الاتجاه \emph{الآخر}: لقد تعلم النموذج القوة ويستمر في تعويضها. وهذا دليل مباشر على أن المتعلَّم نموذج، لا مجرد أن الأمر ”صار ينجح“.

وتُظهر التجارب أمرًا أدق: تتعلم عمليتان في آن واحد، سريعة تتعلم بسرعة وتنسى بسرعة، وبطيئة تتعلم ببطء وتتذكر طويلًا~\cite{Smith2006}. ويُحدَّث التصحيحان بعد كل حركة، عند كل حدث:
\begin{empheq}[box=\keybox]{equation}
e = p - (x_f + x_s), \qquad x_f \leftarrow A_f\,x_f + B_f\,e, \qquad x_s \leftarrow A_s\,x_s + B_s\,e ,
\label{eq:tworate}
\end{empheq}
حيث $p$ الاضطراب، و$x_f$ و$x_s$ التصحيحان المتعلَّمان للعمليتين، و$A$ مقدار ما يُحفظ من التصحيح حتى الحركة التالية، و$B$ مدى قوة تعلمه من الخطأ. للعملية السريعة $B$ كبير و$A$ أصغر من الواحد بوضوح؛ وللبطيئة العكس.

\begin{figure}[t]
\centering
\begin{tikzpicture}[>=Stealth,x=0.034cm,y=1.9cm]
\draw[->,gray!70!black] (0,0) -- (355,0) node[right,font=\small,black] {الحركة};
\draw[->,gray!70!black] (0,-1.15) -- (0,1.25);
\foreach \v in {-1,1}{\draw[gray!60!black] (-3,\v) -- (3,\v); \node[left,font=\scriptsize] at (0,\v) {$\v$};}
\foreach \n in {100,200,300}{\draw[gray!60!black] (\n,-0.04) -- (\n,0.04); \node[below,font=\scriptsize] at (\n,0) {\n};}
\draw[thin,gray!70!black] plot file {figures/tworate_p.dat};
\draw[thick,orange!85!black] plot file {figures/tworate_fast.dat};
\draw[thick,blue!60!black] plot file {figures/tworate_slow.dat};
\draw[very thick,black] plot file {figures/tworate_net.dat};
\node[font=\scriptsize,gray!50!black] at (120,1.12) {الاضطراب $p$};
\node[font=\scriptsize,black,anchor=west] at (238,0.52) {المجموع: عودة};
\node[font=\scriptsize,blue!60!black,anchor=west] at (150,0.39) {البطيئة $x_s$};
\node[font=\scriptsize,orange!85!black,anchor=west] at (60,0.08) {السريعة $x_f$};
\draw[densely dashed,gray!60!black] (226,-1.15) -- (226,1.25);
\node[font=\scriptsize,align=center,anchor=south west] at (228,-1.1) {لا أخطاء\\بعد الآن};
\end{tikzpicture}
\caption{عمليتا التعلم بحسب النموذج~\eqref{eq:tworate}، $A_f=0.92$، $B_f=0.1$، $A_s=0.996$، $B_s=0.02$. مئتا حركة مع الاضطراب $p=1$، ثم ست مع الاضطراب المعاكس $p=-1$، حتى يعود المجموع إلى الصفر. بعد الخط المتقطع يتوقف وصول الأخطاء، فيعود المجموع بنفسه إلى التصحيح القديم: العملية السريعة نسيت الاضطراب المعاكس، والبطيئة تتذكر الأصلي.}
\label{fig:tworate}
\end{figure}

يفسر النموذج~\eqref{eq:tworate} تجربة يصعب فهمها بدونه (الشكل~\ref{fig:tworate}). في حسابنا يبلغ مجموع التصحيحين بعد مئتي حركة مع الاضطراب $0.84$، وتحمل العملية البطيئة معظمه، $0.63$. ثم تعيد ست حركات مع الاضطراب المعاكس المجموع إلى الصفر: السريعة ذهبت إلى السالب، $-0.53$، والبطيئة لا تزال موجبة، $0.46$. فإذا توقف الإبلاغ عن الأخطاء الآن، نسيت العملية السريعة تصحيحها خلال عشرات الحركات، والبطيئة خلال مدة أطول بكثير، فيرتفع المجموع \emph{بنفسه} إلى $0.37$ خلال أربعين حركة: تعود المهارة القديمة بلا أي تدريب. وهذه الاستعادة التلقائية مقيسة لدى البشر؛ أما النموذج ذو العملية الواحدة فلا يستطيع إعطاءها، إذ يخمد بلا أخطاء إلى الصفر ببساطة.

لماذا عمليتان؟ يعطي المرشح الأمثل من الفصل~\ref{sec:acet} جوابًا هندسيًا معقولًا. يتغير الجسم لأسباب مختلفة وبسرعات مختلفة: تتعب العضلات في دقائق، وتنمو وتضعف في أشهر. وإذا كانت الاضطرابات سريعة وبطيئة، فأفضل تقدير يتألف من مرشحين بثابتي زمن مختلفين، يلتقط كل منهما اضطرابه~\cite{Kording2007}. العملية السريعة تصحيح مؤقت، ”اليد تعبت“؛ والبطيئة ”اليد صارت أخرى“. هذا لا يزال فرضية، لكنه يفسر لماذا تضيع المهارة المتعلَّمة في يوم جزئيًا في اليوم التالي، ولماذا تُتعلَّم أسرع في المرة الثانية.

\section{الخلاصة}

\begin{table}[t]
\centering
\footnotesize
\setlength{\tabcolsep}{4pt}
\renewcommand{\arraystretch}{1.15}
\begin{tabular}{>{\raggedright}p{3.0cm}>{\raggedright}p{4.6cm}>{\raggedright\arraybackslash}p{5.2cm}}
\toprule
ماذا تفعل الدارة & كيف & ما هو معروف \\
\midrule
تتعلم بمعلّم & سلك خطأ إلى كل مخرج، القاعدة~\eqref{eq:lms} & بنية الدارة والإضعاف عند التزامن مقيسان؛ والتعلم الموجَّه فرضية رائدة \\
توسّع الدخل & سبعون مليار عصبون \nt{GLUT} & عدد العصبونات مقيس؛ ودور التوسيع نظرية \\
تراعي تأخير الخطأ & أثر مضبوط على التأخير & مقيس \\
تبني نموذجًا داخليًا & مرشح تكيّفي~\eqref{eq:adaptivefilter} & فرضية راسخة الأسس \\
تتعلم بوتيرتين & عمليتان سريعة وبطيئة~\eqref{eq:tworate} & الأثر اللاحق والاستعادة التلقائية مقيسان؛ والعمليتان نموذج \\
\bottomrule
\end{tabular}
\caption{كيف تتعلم الآلة الحركة.}
\label{tab:motorlearning}
\end{table}

صار للآلة الآن ثلاثة معلّمين (الجدول~\ref{tab:motorlearning}). هيب يضغط ما هو متكرر؛ و\nt{DOPA} يعلّم بعدد واحد اختيار المجدي؛ وسلك الخطأ يعلّم الدقة، ويعلّمها بسرعة، لأن كل مخرج يصله خطؤه الخاص. ويبدو أن الدماغ يستعمل الثلاثة، كلًّا حيث يكون أرخص: إشارة البث للقرارات القليلة، والأسلاك المستقلة للضبط الدقيق للجسم، حيث تُبلغ المستشعرات نفسها عن الجواب الصحيح.

\section{تمارين}

\begin{exercise}[\lvlA]
\label{ex:15:1}
\emph{سعة الدارة ذات سلك الخطأ.} في الدارة $3\cdot10^7$ عصبون مخرج لكل منها $10^5$ مدخل، ولكل منها سلك خطأ خاص (نبضة واحدة في الثانية). العصبون ذو $n$ مدخلًا يتعلم أي تصنيف لـ$\approx2n$ متجهًا~\cite{Cover1965}. (أ)~كم ارتباطًا تخزن الدارة؟ (ب)~قدّر بحسب~\eqref{eq:spikeentropy} عند $\Delta t=10\ms$ أقصر زمن لملء سعة العصبون.
\end{exercise}

\begin{exercise}[\lvlB]
\label{ex:15:2}
\emph{سرعة قاعدة المربعات الصغرى.} أنماط القاعدة~\eqref{eq:lms} تخمد بسرعات $\eta\lambda_k(C)$. لخمسة مرشحات~\eqref{eq:adaptivefilter} على ضجيج أبيض $C_{jk}=2\sqrt{\tau_j\tau_k}/(\tau_j+\tau_k)$. أوجد نسبة أسرع السرعات إلى أبطئها إذا تزايدت $\tau_j$ بعامل $2$ ($10$ إلى $160\ms$) وبعامل $4$.
\end{exercise}

\begin{exercise}[\lvlB]
\label{ex:15:3}
\emph{تحليل العمليتين.} اكتب النموذج~\eqref{eq:tworate} بمعاملات الشكل~\ref{fig:tworate} على الصورة $\mathbf x_{n+1}=M\mathbf x_n+\mathbf b\,p$. (أ)~أوجد من القيم الذاتية لـ$M$ الثوابت الزمنية مقيسة بالحركات. (ب)~أوجد قيمتي $x_f$ و$x_s$ المستقرتين ومجموعهما عند $p=1$ ثابت.
\end{exercise}

\begin{exercise}[\lvlB]
\label{ex:15:4}
\emph{نمذجة: المرة الثانية أسرع.} بالنموذج~\eqref{eq:tworate} بمعاملات من \foreignlanguage{english}{\texttt{models/\allowbreak motor\_learning.py}} أوجد عدد الحركات حتى $x_f+x_s\ge0.8$ عند $p=1$: (أ)~لآلة غير مدرَّبة؛ (ب)~بعد $200$ حركة مع $p=1$ واضطراب معاكس حتى صفر المجموع، كما في الشكل~\ref{fig:tworate}. هل يمكن أن يتسارع هكذا نموذج بعملية واحدة؟
\end{exercise}

\begin{exercise}[\lvlB]
\label{ex:15:5}
\emph{منظِّم بنموذج داخلي.} الجسم المتحكَّم فيه $\dd x/\dd t=ku$ يُرى بتأخير $d$؛ والمنظِّم $u=-G\hat x$ يتنبأ بـ$\hat x(t)=x(t-d)+\hat k\int_{t-d}^{t}u\,\dd s$. (أ)~عند $\hat k=k=1$ و$d=150\ms$ أوجد $G$ لثابت زمن $20\ms$ وقارنه بالحد~\eqref{eq:delaystable}. (ب)~هل الدارة مستقرة عند $\hat k=0.4k$؟
\end{exercise}

\begin{exercise}[\lvlB]
\label{ex:15:6}
\emph{المرشح التكيّفي يتعلم العضلة.} الجسم المتحكَّم فيه نفضة~\eqref{eq:twitch} مساحتها واحد، $T_c=50\ms$؛ والمرشحات~\eqref{eq:adaptivefilter} دارات \foreignlanguage{english}{RC} بـ$\tau_j=12.5$، $25$، $50$، $100$، $200\ms$؛ والدخل عدد طبيعي التوزيع جديد كل $10\ms$. في كم ثانية تخفض القاعدة~\eqref{eq:lms} الخطأ إلى $0.5\%$ عند $\eta=50$ و$200$؟ ولماذا تبقى الأوزان بعيدة عن الأفضل حتى بعد $300\,\text{s}$؟
\end{exercise}

\begin{exercise}[\lvlC]
\label{ex:15:7}
\emph{العمليتان مرشحًا أمثل.} الاضطراب مجموع عمليتين $p^{f,s}_{n+1}=A_{f,s}\,p^{f,s}_n+w^{f,s}_n$ بتباينَي ضجيج $q_f$ و$q_s$، والخطأ يُرصد بضجيج تباينه $R$؛ ومرشح كالمان يعطي~\eqref{eq:tworate}. عند أي $q_f/R$ و$q_s/R$ ينتج من $A_f=0.92$ و$A_s=0.996$ أن $B_f=0.1$ و$B_s=0.02$؟ وكيف يتغير $B_f$ و$B_s$ إذا ضوعف $q_s$ أربع مرات؟
\end{exercise}
